\documentclass[11pt]{article}

\usepackage[a4paper,margin=1in]{geometry}
\usepackage{amsmath,amssymb,amsthm,booktabs,microtype}
\usepackage{url}

\pdfinfo{
  /Title (Exact Solution and Optimal Face of a Four-Dimensional Boolean-Quadric Semidefinite Relaxation)
  /Author (Xiaojian Zeng)
}

\newtheorem{theorem}{Theorem}
\newtheorem{lemma}{Lemma}
\newcommand{\R}{\mathbb{R}}
\newcommand{\B}{\mathcal{B}}
\newcommand{\F}{\mathcal{F}}
\newcommand{\one}{\mathbf{1}}
\newcommand{\ip}[2]{\left\langle #1,#2\right\rangle}

\title{Exact Solution and Optimal Face of a Four-Dimensional\\
Boolean-Quadric Semidefinite Relaxation}
\author{Xiaojian Zeng\\
\small PBC School of Finance, Tsinghua University, Beijing 100190, China\\
\small \texttt{xiaojianzeng@foxmail.com}}
\date{}

\begin{document}
\maketitle

\begin{abstract}
We solve exactly a specific four-dimensional semidefinite relaxation introduced
by Zhang and Wang: their Boolean-quadric-polytope relaxation (9) for the
box-constrained quadratic instance (5) with $a=b=1/2$.  If $\alpha$ is the
unique root in $(-1/24,-1/25)$ of
$32\alpha^3+48\alpha^2+26\alpha+1=0$, then the optimum is
\[
 -\frac{32\alpha^2+16\alpha+1}{36}
 =-0.0108411477500751\ldots.
\]
We give matching exact primal and dual certificates.  Strict dual slacks and
semidefinite complementarity also show that the complete optimal face,
including the auxiliary probability vector used to represent the Boolean
quadric polytope, is a singleton.  This replaces the previously available
$-1/100$ feasible value and numerical value near $-0.0108$ by an exact
certificate and a full face classification.
\end{abstract}

\section{The finite conic program}

Boolean quadric polytopes provide a standard lifted representation of binary
quadratic products~\cite{Padberg1989}; their connection with box-constrained
nonconvex quadratic programming is developed in~\cite{BurerLetchford2009}.
Zhang and Wang~\cite{ZhangWang2026} consider
\begin{equation}
Q=\begin{pmatrix}
1/4&-1/2&0&0\\
-1/2&1&-1&0\\
0&-1&1&-1/2\\
0&0&-1/2&1/4
\end{pmatrix},
\qquad c=(1/2,1,0,0)^\top.                                       \tag{1}
\end{equation}
For $x\in\R^4$, $X\in\mathbb{S}^4$, and a probability vector
$p=(p_z)_{z\in\{0,1\}^4}$, define
\[
Y(x,X)=\begin{pmatrix}1&x^\top\\x&X\end{pmatrix}.
\]
The program addressed here is
\begin{equation}
\begin{aligned}
v_*={}&\min\ \ip{Q}{X}+c^\top x,\\
&Y(x,X)\succeq0,\qquad p_z\geq0,\qquad \sum_z p_z=1,\\
&x_i=\sum_z p_z z_i,\qquad
X_{ij}=\sum_z p_z z_i z_j\quad(1\leq i<j\leq4).
\end{aligned}                                                     \tag{2}
\end{equation}
Thus the full off-diagonal projection belongs to
\[
\B_4=\operatorname{conv}\{(z,(z_i z_j)_{i<j}):z\in\{0,1\}^4\},
\]
while the diagonal entries $X_{ii}$ are constrained through $Y\succeq0$
rather than set equal to $x_i$.  This is the full formulation of relaxation
(9), not the relaxation containing only the displayed McCormick inequalities.

The source proves that the underlying box-constrained quadratic program has
value zero, gives a feasible point of (2) with value $-1/100$, and reports an
uncertified numerical optimum near $-0.0108$~\cite[Propositions 1--3 and
Appendix A]{ZhangWang2026}.  The result below closes precisely that finite
instance.  No claim about other parameter values or about priority beyond this
comparison is intended.

\section{Exact certificate}

Put
\[
P(A)=32A^3+48A^2+26A+1.
\]
Since $P(-1/24)=-1/432$, $P(-1/25)=543/15625$, and $P'$ is positive on
$[-1/24,-1/25]$, there is a unique
\begin{equation}
\alpha\in(-1/24,-1/25)\quad\text{with}\quad P(\alpha)=0.          \tag{3}
\end{equation}
Define
\begin{equation}
t=\frac{32\alpha^2+16\alpha+1}{36},\qquad
v=\frac{8\alpha^2+8\alpha+1}{2},\qquad w=-2\alpha-1.             \tag{4}
\end{equation}

All unlisted probabilities in the following display are zero:
\begin{equation}
\begin{array}{c@{\qquad}l}
\toprule
z & p_z\\
\midrule
0000,1111 & \displaystyle \frac{4(5144\alpha^2+3256\alpha+1711)}{35073}\\[3pt]
0001,0111 & \displaystyle -\frac{2(7552\alpha^2+3362\alpha-2779)}{35073}\\[3pt]
0011 & \displaystyle -\frac{8(1046\alpha^2+1153\alpha-1193)}{35073}\\[3pt]
1010 & \displaystyle -\frac{2576\alpha^2+3376\alpha-725}{35073}\\
\bottomrule
\end{array}                                                       \tag{5}
\end{equation}
Use these probabilities to define $x_i=\sum_z p_z z_i$ and
$X_{ij}=\sum_z p_z z_i z_j$ for $i<j$, and set
\begin{align}
X_{11}&=\frac{8368\alpha^2+9224\alpha+2147}{11691},\nonumber\\
X_{22}&=-\frac{4(1204\alpha^2+53\alpha-2245)}{35073},\nonumber\\
X_{33}&=-\frac{15760\alpha^2+12812\alpha-19249}{35073},\nonumber\\
X_{44}&=-\frac{8(454\alpha^2-349\alpha-1099)}{11691}.            \tag{6}
\end{align}
The six ordered weights on
$(0000,0001,0011,0111,1010,1111)$ are approximately
\[
(0.1807155517,0.1656936580,0.2826366501,
  0.1656936580,0.02454493068,0.1807155517).
\]

With the scalar simplex multiplier set to $\tau=0$, the dual matrix is
\begin{equation}
S=\begin{pmatrix}
t&\frac14-\alpha-v&-\frac12-v-w&\alpha&0\\
\frac14-\alpha-v&\frac14&-\frac12-\alpha&v&2\alpha\\
-\frac12-v-w&-\frac12-\alpha&1&w&v\\
\alpha&v&w&1&-\frac12-\alpha\\
0&2\alpha&v&-\frac12-\alpha&\frac14
\end{pmatrix}.                                                    \tag{7}
\end{equation}

\begin{theorem}\label{thm:main}
The optimum of \emph{(2)} is
\[
v_*=-t=-\frac{32\alpha^2+16\alpha+1}{36}.
\]
Equivalently, $v_*$ is the root near $-0.0108411477500751$ of
\begin{equation}
576V^3-16V^2+92V+1=0.                                             \tag{8}
\end{equation}
The complete optimal face in $(x,X,p)$ is the single point
\emph{(5)}--\emph{(6)} with the moment definitions above.  This point is
fixed by reverse-complement symmetry.
\end{theorem}

\begin{proof}
We first prove dual feasibility.  Let $B$ be the lower-right $4\times4$
block of $S$.  In orthonormal bases adapted to reversal of the four
coordinates, its even and odd blocks are
\begin{equation}
B_+=\begin{pmatrix}
\frac14+2\alpha&-\frac12-\alpha+v\\
-\frac12-\alpha+v&1+w
\end{pmatrix},\qquad
B_-=\begin{pmatrix}
\frac14-2\alpha&-\frac12-\alpha-v\\
-\frac12-\alpha-v&1-w
\end{pmatrix}.                                                     \tag{9}
\end{equation}
Reduction modulo $P(\alpha)=0$ gives
\[
\det B_+=\det B_-=0.
\]
All four diagonal entries in (9) are positive on the isolating interval
(3), so both blocks are positive semidefinite of rank one.  If
\[
k_1=\tfrac14-\alpha-v,\qquad k_2=-\tfrac12-v-w,
\]
then the top row of $S$ lies in the ranges of these two blocks because
\begin{align*}
(\tfrac14+2\alpha)(k_2+\alpha)-(-\tfrac12-\alpha+v)k_1&=0,\\
(\tfrac14-2\alpha)(k_2-\alpha)-(-\tfrac12-\alpha-v)k_1&=0.
\end{align*}
Its Schur complement is zero:
\begin{equation}
t-\frac{k_1^2}{2(\frac14+2\alpha)}
 -\frac{k_1^2}{2(\frac14-2\alpha)}=0.                              \tag{10}
\end{equation}
The identities again hold exactly modulo $P$.  Hence $S\succeq0$ and
$\operatorname{rank}S=2$.

For $z\in\{0,1\}^4$, let $F_z\in\mathbb{S}^5$ have top row and column
$z$, lower off-diagonal entries $(F_z)_{ij}=z_i z_j$, and zero diagonal.
Writing $d_i=X_{ii}$, every feasible lifted matrix has the form
\begin{equation}
Y=E_{00}+\sum_z p_zF_z+\sum_{i=1}^4d_iE_{ii}.                      \tag{11}
\end{equation}
Set
\[
h_z=c^\top z+2\sum_{i<j}Q_{ij}z_i z_j,\qquad
r_z=h_z-\ip{S}{F_z}.
\]
The sixteen exact slacks are zero precisely for
\begin{equation}
0000,\ 0001,\ 0011,\ 0111,\ 1010,\ 1111,                         \tag{12}
\end{equation}
and the remaining values, grouped when equal, are
\begin{equation}
\begin{array}{c@{\quad}c}
\toprule
\text{atoms} & \text{slack}\\
\midrule
0010,1011 & -2\alpha\\
0101 & -4\alpha\\
0100,1101 & 8\alpha^2+4\alpha+1\\
0110,1001 & (2\alpha+1)(4\alpha+1)\\
1000,1110 & 8\alpha^2+10\alpha+1\\
1100 & 2(8\alpha^2+8\alpha+1)\\
\bottomrule
\end{array}                                                         \tag{13}
\end{equation}
Every expression in (13) is strictly positive on (3).  Also
$S_{ii}=Q_{ii}$ for $i=1,\ldots,4$.  It follows from (11) that every
feasible point satisfies the exact weak-duality identity
\begin{equation}
\ip{Q}{X}+c^\top x
=\ip{S}{Y}-t+\sum_z p_zr_z\geq-t.                                 \tag{14}
\end{equation}

The probabilities (5) are strictly positive and sum to one; these facts
follow by reducing the sum modulo $P$ and applying rational interval
arithmetic on (3).  Direct multiplication in $\mathbb{Q}[\alpha]/(P)$ gives
\begin{equation}
SY=0,\qquad \ip{Q}{X}+c^\top x=-t.                                \tag{15}
\end{equation}
It remains only to justify $Y\succeq0$.  The leading principal minors of
its submatrix on indices $\{0,1,2\}$, after the initial $1$, are
\begin{equation}
\frac{4(4064\alpha^2+5146\alpha+1171)}{35073},\qquad
-\frac{8(313756\alpha^2+607256\alpha-428197)}{1230115329},          \tag{16}
\end{equation}
and both are positive on (3).  Thus this $3\times3$ principal submatrix
is positive definite.  Since $S$ has rank two and $SY=0$, we have
$\operatorname{rank}Y\leq3$; (16) forces equality.  The Schur complement
of the positive definite block therefore has rank zero, proving
$Y\succeq0$.  Equations (14)--(16) give matching exact dual and primal
certificates and establish the value.

For completeness, consider an arbitrary optimal point.  Equality in (14)
and the strict slacks (13) force $p_z=0$ outside (12).  Moreover,
$\ip{S}{Y}=0$ for two positive semidefinite matrices, hence $SY=0$.
Use the six probabilities in (12) and the four diagonal variables as the
ten unknowns.  A square subsystem of normalization and $SY=0$---the five
entries of row zero and columns $0,1,2,4$ of row one---has determinant
\begin{equation}
\frac{5775280\alpha^2+8296838\alpha+334859}{331776}.               \tag{17}
\end{equation}
The numerator in (17) is coprime to $P$, so it cannot vanish at $\alpha$.
The variables (5)--(6) satisfy the full system, and (17) makes them its
unique solution.  The moment equations then determine $x$ and every
off-diagonal entry of $X$ uniquely.  Thus the entire optimal face is the
claimed singleton.

Finally, reverse-complement sends
$0000\leftrightarrow1111$, $0001\leftrightarrow0111$, and fixes $0011$
and $1010$.  The equal weights in (5), together with direct substitution
in (6), shows that the singleton is fixed.  Eliminating $\alpha$ between
(3) and $V=-t$ gives (8).
\end{proof}

\section{Verification and scope}

All displayed equalities can be checked using degree-two representatives in
$\mathbb{Q}[\alpha]/(P)$, and every strict sign can be checked by rational
interval arithmetic on $[-1/24,-1/25]$.  The accompanying
\texttt{exact\_certificate.py} carries out these reductions, enumerates all
sixteen atom slacks, reconstructs the primal moments, and checks the rank-ten
uniqueness system.  Exact rational or algebraic sign certification with
inspectable verification programs has also been used in finite-dimensional
matrix-positivity problems~\cite{ZengEtAl2026}; the present certificate is
independent and specific to (2).

The result concerns only the fixed $n=4$, $a=b=1/2$ instance of (2).  It
does not strengthen the asymptotic relaxation-gap theorems
of~\cite{ZhangWang2026}, and it does not classify optima for other parameters
or for relaxations with a different lifted feasible region.

\paragraph{Computational and AI assistance.}
The algebraic certificate was discovered and checked with computer algebra.
OpenAI Codex assisted with symbolic exploration, exact-certificate code, and
manuscript preparation.  The released Python/SymPy verifier uses exact rational
polynomial reduction and rational interval arithmetic; floating-point values in
this article are descriptive only and are not used in the proof.

\small
\bibliographystyle{plain}
\bibliography{references}

\end{document}
